\section{Closure by an ansatz for the distribution}
\label{sec:moment}

When no small parameter is available, or when a model must stay valid away
from equilibrium, the alternative is to retain more moments and assume a
shape for $f$ with as many free parameters as retained moments. The missing
moment is then computed from that shape. These closures give evolution
equations for the fluxes, so the fluxes have finite relaxation times and
signals have finite speeds.

\subsection{The linear model problem in a Hermite basis}
\label{sec:hermite}

One linear problem is used repeatedly below. Take one space and one
velocity dimension, units with $n_0=m=\vt=1$, the Maxwellian
$M(v)=(2\pi)^{-1/2}e^{-v^2/2}$, and write $f=M(1+h)$. The linearized
kinetic equation with acceleration $a(x,t)$ is
\begin{equation}
  \partial_th+v\,\partial_xh+\Lambda h=a\,v .
  \label{eq:linkin}
\end{equation}
Expand $h=\sum_{n\ge0}a_n(x,t)\,\psi_n(v)$ in the orthonormal Hermite
functions $\psi_n=\He_n/\sqrt{n!}$, where $\He_0=1$, $\He_1=v$,
$\He_2=v^2-1$, $\He_3=v^3-3v$. The recurrence
$v\,\psi_n=\sqrt{n+1}\,\psi_{n+1}+\sqrt n\,\psi_{n-1}$ turns
\eqref{eq:linkin} into
\begin{equation}
  \partial_ta_n+\partial_x\bigl(\sqrt{n+1}\,a_{n+1}+\sqrt n\,a_{n-1}\bigr)
  +(\Lambda a)_n=a\,\delta_{n1} .
  \label{eq:hermitehier}
\end{equation}
The coefficients are the fluid moments: $a_0$ is the density perturbation,
$a_1$ the velocity, $\sqrt2\,a_2$ the temperature, $\sqrt6\,a_3$ the heat
flux (third central moment), $\sqrt{24}\,a_4$ the deviation of the fourth
moment from its Maxwellian value. For BGK, $\Lambda$ is diagonal with
entries $0,0,0,\nu,\nu,\dots$; for the Dougherty (Lenard--Bernstein)
operator \cite{LenardBernstein58,Dougherty64} the entries are
$0,0,0,3\nu,4\nu,\dots$.

Streaming is a tridiagonal (Jacobi) matrix in this basis. It couples each
coefficient to its two neighbours and nothing else, and all closure
questions for the linear problem are questions about what to do at the last
retained row.

\subsection{Truncation}

The simplest choice is $a_{N+1}=0$. For a linear problem this is Grad's
closure. Its content is exact and worth knowing.

\begin{example}[Truncation is a discrete-velocity model]
With $a_{N+1}=0$ the collisionless system is
$\partial_t\mathbf a+J_N\,\partial_x\mathbf a=0$, where $J_N$ is the
$(N+1)\times(N+1)$ Jacobi matrix with off-diagonal entries
$\sqrt1,\dots,\sqrt N$. Its eigenvalues are the roots of $\He_{N+1}$, the
Gauss--Hermite quadrature nodes. (Checked for $N=8$: the eigenvalues and
the nodes agree to $3\times10^{-15}$, and the largest is $4.51$.) The
truncated moment system is therefore equivalent to a kinetic equation with
$N+1$ discrete velocities and Gauss--Hermite weights.
\end{example}

Two consequences follow. The truncated system is hyperbolic with real wave
speeds, so it is well posed. And without collisions it cannot damp: a
density perturbation at wavenumber $k$ is a finite sum of oscillations
$e^{-ikv_jt}$, which decays for a while and then returns. The recurrence
time is of order $\sqrt N/(|k|\vt)$. Landau damping
(Section~\ref{sec:collisionless}) is lost, however many moments are kept,
unless something absorbs the flow of perturbation energy to high $n$. That
absorber is either collisions, which act strongly at high $n$, or a
closure that models the outflow.

\subsection{Grad's method}

Grad \cite{Grad49} expands the nonlinear distribution as a local Maxwellian
times a polynomial,
\[
  f=f_M(n,\bu,T)\Bigl[1+\frac{\boldsymbol\pi:\bc\bc}{2p}
     +\frac{2}{5}\,\frac{\bq\cdot\bc}{p\,\vt}
       \Bigl(\frac{|\bc|^2}{2}-\frac52\Bigr)\Bigr],
\]
which has thirteen parameters $(n,\bu,T,\boldsymbol\pi,\bq)$, and takes
moments of the Boltzmann equation to get thirteen evolution equations. The
unclosed higher moments are evaluated with this $f$. The standard counts
are 5 (Euler), 10 (Gaussian), 13, 14, 20, 26 and 35 moments.

Grad's equations contain Navier--Stokes--Fourier as the limit in which
$\boldsymbol\pi$ and $\bq$ relax quickly. They have two known defects. The
ansatz is not positive, and the equations lose hyperbolicity at moderate
departures from equilibrium, which appears as unphysical sub-shocks.
Regularized (R13) and globally hyperbolic variants address these
\cite{StruchtrupTorrilhon03,Struchtrup05}.

\subsection{Maximum entropy}

Levermore \cite{Levermore96} chooses the distribution that maximizes the
entropy $-\int(f\ln f-f)\,\dd v$ subject to the retained moments
$\int m_i(\bv)f\,\dd v=\rho_i$. The maximizer is
\begin{equation}
  f=\exp\Bigl(\sum_i\alpha_i\,m_i(\bv)\Bigr),
\end{equation}
with multipliers $\alpha_i$ determined by the moments. The resulting system
has three structural properties by construction: it is symmetric hyperbolic
in the variables $\alpha$, it dissipates an entropy, and the ansatz is
positive, so its moments are realizable.

The costs are also structural. The map from moments to multipliers has a
closed form only for the 5-moment (Euler) and 10-moment (Gaussian) systems,
so each flux evaluation is an optimization problem. And beyond ten moments
the set of moments for which the problem is solvable is not what one would
hope: Junk showed that for the simplest higher system equilibrium lies on
the boundary of that set, and the fluxes are singular there \cite{Junk98}.
A closed-form example in another setting, the M1 closure for radiation, is
worked in Section~\ref{sec:rad}.

\subsection{Realizability and hyperbolicity}
\label{sec:realizability}

A vector of moments is \emph{realizable} if it comes from some non-negative
distribution. In one velocity dimension this is the Hamburger moment
problem: $(M_0,\dots,M_{2r})$ is realizable when the Hankel matrices
$(M_{i+j})_{i,j=0}^{r}$ are positive semidefinite. The first conditions are
$M_0\ge0$, $M_0M_2-M_1^2\ge0$ (non-negative temperature), and a lower bound
on the fourth moment in terms of the lower ones. A closure can leave the
realizable set in two ways: its formula for the unclosed moment can be
inconsistent with the retained ones, or the closed dynamics can carry
realizable data out of the set.

\emph{Hyperbolicity} is the requirement that the flux Jacobian of the
closed first-order system have real eigenvalues and a full set of
eigenvectors. Without it the initial-value problem is ill posed, and no
amount of accuracy in the closure term compensates. For learned closures
this is the property most often lost when a network is fitted to the
unclosed moment directly (Section~\ref{sec:learned}).

\subsection{Other ansatz closures}

\begin{itemize}
  \item \textbf{Gaussian (10-moment).} Retain the full pressure tensor and
    set the heat flux to zero. Used in plasma modelling where pressure
    anisotropy matters; the heat-flux closure added on top is often a
    relaxation term built on the Hammett--Perkins idea \cite{Wang15}.
  \item \textbf{Spectral truncations.} $P_N$ in radiation transport is
    truncation of a spherical-harmonic expansion; it is linear and can
    produce negative densities.
  \item \textbf{Quadrature-based methods.} Represent $f$ as a sum of delta
    functions (or kernels) whose weights and nodes are recovered from the
    retained moments \cite{McGraw97}. Realizable by construction; the
    inversion is ill conditioned at high order.
\end{itemize}

\begin{center}
\renewcommand{\arraystretch}{1.25}
\begin{tabularx}{\textwidth}{@{}P{2.9cm}P{3.5cm}YY@{}}
\toprule
Closure & Ansatz & Guarantees & Known failures \\
\midrule
Truncation / Grad & Maxwellian times polynomial
  & linear problem hyperbolic; recovers NSF
  & not positive; nonlinear loss of hyperbolicity; no collisionless damping \\
Maximum entropy & exponential of polynomial
  & hyperbolic, entropy, realizable
  & costly; singular fluxes beyond 10 moments \\
Quadrature & sum of deltas
  & realizable
  & ill-conditioned inversion; weakly hyperbolic variants \\
Gaussian & anisotropic Maxwellian
  & hyperbolic, realizable, closed form
  & no heat flux \\
\bottomrule
\end{tabularx}
\end{center}
